\subsection{Stored bytes and allocated bytes}

The eight-bit-residual configuration has an int8-scale resident
representation. For one-byte primary subspace codes, its row payload is
\begin{equation}
 S_{\rm rows}=N\left(M+\left\lceil\frac{dB_r}{8}\right\rceil+8\right),
 \qquad S_{\rm offsets}=4(n_{\rm list}+1).
 \label{eq:4-row-storage}
\end{equation}
The extra eight bytes store a four-byte ID and FP32 correction. At
$M=d/8$, $B_r=8$, codes occupy $9d/8$ bytes: 12.5\% more than an SQ8
vector's $d$ bytes. A graph-free hierarchy is therefore a design choice,
not evidence of a unique compression advantage over int8 alternatives.

\begin{table}[tbp]
\centering
\caption[Per-vector storage model]{Per-vector bytes. JHQ includes codes, ID and correction; SQ8 is
vector codes alone; CAGRA int8 models vector codes plus 32 four-byte
neighbours. Offsets, model parameters, workspace and other runtime
allocations are excluded. These components are not complete process totals.}
\label{tab:4-bytes}
\footnotesize\setlength{\tabcolsep}{4pt}
\begin{tabular}{@{}rrrrrrrr@{}}
\toprule
$d$ & Primary & Residual & ID + corr. & JHQ & SQ8 & \shortstack{CAGRA\\int8} & FP32 \\
\midrule
768 & 96 & 768 & 8 & 872 & 768 & 896 & 3,072 \\
1,024 & 128 & 1,024 & 8 & 1,160 & 1,024 & 1,152 & 4,096 \\
1,536 & 192 & 1,536 & 8 & 1,736 & 1,536 & 1,664 & 6,144 \\
3,072 & 384 & 3,072 & 8 & 3,464 & 3,072 & 3,200 & 12,288 \\
\bottomrule
\end{tabular}
\end{table}

The CAGRA comparison uses $d+4d_g$ bytes at graph degree $d_g=32$,
consistent with the harness's 32-bit neighbour IDs. On Stella, JHQ's row
payload is 19.20\,GiB, while a successful degree-32 CAGRA int8 record
reports 19,656\,MiB (19.20\,GiB) of index allocation. The near equality is
useful context despite their different accounting boundaries. Chapter 7
reports both int8 and fp32 allocations alongside search quality and cost.
JHQ's distinctions are graph-free construction and a query-adjustable
residual-access budget; lowering $\alpha$ leaves resident residual codes
intact.

The basic FP32 model arrays require $4d^2$ bytes for the transform,
$4MKD_s$ for the materialised primary codebook, $4MK_r$ for residual
codebooks and $4n_{\rm list}(d+1)$ for coarse centroids and norms.
Derived centroid copies and encoder levels add allocations in the
selected path. Query workspace includes two FP32 query matrices
$[b_{\rm cap},d]$, centroid scores $[b_{\rm cap},n_{\rm list}]$, LUTs,
probes and survivor buffers. The optional residual LUT would add
$4b_{\rm cap}dK_r$ bytes; default refinement computes residual scores
directly. Group scratch follows \cref{eq:5-group-workspace} and can retain
a larger capacity than the current scan uses: recorded 128-probe samples
reserve about 93.8\,MiB even on the query-centric path.

\begin{table}[tbp]
\centering
\caption[Construction and query object lifetimes]{Object lifetimes. Sizes describe named arrays, rather than all
allocations of a CUDA process. Construction phases are expanded in
\cref{sec:4-5}.}
\label{tab:4-lifetimes}
\footnotesize\setlength{\tabcolsep}{4pt}
\begin{tabularx}{\linewidth}{@{}>{\raggedright\arraybackslash}p{2.0cm}>{\raggedright\arraybackslash}p{1.7cm}>{\raggedright\arraybackslash}p{1.5cm}*{2}{>{\raggedright\arraybackslash}X}@{}}
\toprule
Object & Producer & Location & Lifetime & Bytes / scale \\
\midrule
Input / samples & Loader / harness & Host & Build input; external ownership & $4Nd$ raw; sampled subsets \\
Transform / models & QR / training & Host, GPU & Model and index lifetime & $4d^2+4MKD_s+4MK_r+4n_{\rm list}(d+1)$ basic GPU arrays \\
Trained cache & Training & Host file & Across processes; rebuild still needed & Serialised model state \\
Input staging & Add pipeline & Pinned host, GPU & During add & $4n_ba_{\rm in}d$ host; $8n_ba_{\rm in}d$ GPU floats, plus scratch \\
Assignment products & IVF assignment & GPU & During add; freed before gather & $4n_b n_{\rm list}a_{\rm assign}$; independent row capacity \\
IDs / offsets & Sort / lower bound & GPU & Completed snapshot & $4N+4(n_{\rm list}+1)$ \\
Codes / correction & Encode / reorder & GPU & Completed snapshot & $N(M+b_r+4)$ \\
Order / inverse & Sort / inversion & GPU & Construction only & $4N$ per array; phase dependent \\
Query workspace / graph & Search / capture & GPU, host & Allocations reused; batch contents replaced & Capacity dependent; graph references resident arrays \\
\bottomrule
\end{tabularx}
\end{table}

\subsection{Reuse, recovery and invalidation}

The cache saves the rotation, codebooks and centroids. Its key mixes model
parameters and a strided checksum of the training sample, but omits the
code version and a checksum of every database byte. Reproducing a build
therefore requires model provenance and explicit cache invalidation when
implementation or training data change. The search experiments reuse
trained states; cold-build records construct separate instances. A close
recall result across those instances cannot establish byte-identical
models.

A cache hit still executes constructor QR and library warm-up before the
training call checks the cache, then loads/uploads the model, encodes
rows and builds lists. Stella's 6.67\,s add record is one component of that
recovery path, not a measured process-restart time. The add call includes its initial workspace allocation; schedule
calibration and search warm-up lie outside that isolated add number.

At fixed $k$, changing $\alpha$ changes the survivor budget without
retraining or adding vectors. Search grows insufficient workspace before
capture, and recaptures when the actual batch size, survivor count,
probe count, refinement lanes, scan groups or shared-bound setting
changes. Calibration can bypass capture while trying budgets. Releasing
the index destroys its graph and device allocations; buffers are
index-owned and reused across the sequential calls evaluated here.
